\documentclass[11pt]{article}
\usepackage[margin=1in]{geometry}
\usepackage{amsmath,amssymb,amsthm,booktabs}
\newtheorem{proposition}{Proposition}
\newcommand{\R}{\mathbb{R}}
\title{From $\ell_\infty$ windows to probabilistic coverage:\\ estimating $\sigma$ and a per-image error ellipsoid}
\date{}
\begin{document}
\maketitle

\section{Deterministic set (current method)}
Camera $i$ has rotation $R_i$, centre $C_i$ and focal length $f$, principal point $(c_x,c_y)$; point $k$ is $X_k$.
In camera coordinates $v_{ik}=R_i(X_k-C_i)$, and the predicted pixel is
\begin{equation}
  \hat y_{ik}=\Bigl(f\,\frac{v_{ik,x}}{v_{ik,z}}+c_x,\ \ f\,\frac{v_{ik,y}}{v_{ik,z}}+c_y\Bigr).
\end{equation}
$y_{ik}$ is the SIFT detection of point $k$ in image $i$. With $\theta=(R_i,C_i)_i$, the detection error at the truth is
$e_{ik}=y_{ik}-\hat y_{ik}(\theta^*,X^*)\in\R^2$, and for windows $\epsilon_{ik}$ the feasible set is
\begin{equation}
  \mathcal F(\epsilon)=\bigl\{(\theta,X):\ \|y_{ik}-\hat y_{ik}(\theta,X)\|_\infty\le\epsilon_{ik}\ \ \text{for every detection}\bigr\}.
\end{equation}
With known rotations each window is a pyramid with apex $C_i$: 4 linear rows per detection, plus a depth row.

\begin{proposition}[deterministic coverage]
If $\|e_{ik}\|_\infty\le\epsilon_{ik}$ for every detection, then $(\theta^*,X^*)\in\mathcal F(\epsilon)$ (and conversely,
since every row is evaluated at $X^*$). For the poses alone, $\theta^*$ is in the projection of $\mathcal F(\epsilon)$
iff for every point $k$ there is \emph{some} $X_k$ (not necessarily $X_k^*$) satisfying all of its rows at $\theta^*$.
\end{proposition}

\section{The probabilistic problem}
The errors $e_{ik}$ are random (image noise, detector). We do not fix their distribution; we only use moments that can be
estimated without ground truth (Section~3). The goal is a set $\Theta$ computed without ground truth such that
\begin{equation}
  \Pr(\theta^*\in\Theta)\ \ge\ 1-\alpha .
\end{equation}

\paragraph{Events.} An \emph{event} is a statement about the errors that is either true or false for a given set of
images. Choose an event $\mathsf E$ such that \emph{whenever $\mathsf E$ is true, $\theta^*\in\Theta$} (deterministic
step, as in the proposition). Then $\Pr(\theta^*\in\Theta)\ge\Pr(\mathsf E)$, and it remains to bound $\Pr(\mathsf E)$.
For the $\ell_\infty$ set, $\mathsf E$ is ``every coordinate of every $e_{ik}$ is at most $\epsilon_{ik}$''.
If $\mathsf E=\mathsf E_1\cap\dots\cap\mathsf E_m$ (all of $\mathsf E_1,\dots,\mathsf E_m$ true), the union bound gives
\begin{equation}
  \Pr(\mathsf E)\ \ge\ 1-\sum_{j}\Pr(\mathsf E_j\ \text{false}),
  \label{eq:union}
\end{equation}
with no independence assumption between the $\mathsf E_j$.

\paragraph{Assumptions.} Correct correspondences (gross mismatches removed); known rotations (or a separate rotation bound).

\section{Estimating $\sigma$ without ground truth}
\paragraph{Which error matters.} Part of the errors of point $k$ is consistent with a shifted point $X_k^*+\delta_k$
(viewpoint-independent detector bias). By the proposition this part is absorbed by the free $X_k$ and does not affect
pose coverage. COLMAP's residual absorbs more than that (the detection is part of the fit), so it underestimates.

\paragraph{Leave-one-out residual.} With COLMAP's poses $\hat\theta$, re-triangulate point $k$ without image $i$
(least squares over its other detections, giving $\hat X_{k,-i}$) and measure
$r_{ik}=y_{ik}-\hat y_{ik}(\hat\theta,\hat X_{k,-i})$. Per coordinate, with $M$ detections,
\begin{equation}
  \hat\sigma^2=\frac{1}{2M}\sum_{ik}\|r_{ik}\|^2,\qquad
  \hat\kappa=\frac{1}{2M\hat\sigma^4}\sum_{ik}\bigl(r_{ik,u}^4+r_{ik,v}^4\bigr)\quad(\kappa=3\text{ for a Gaussian}).
\end{equation}

\paragraph{30-view chair spiral} (RMS per coordinate; ground truth only for evaluation):
\begin{center}
\begin{tabular}{lccc}
\toprule
 & COLMAP residual & leave-one-out (no GT) & error to mesh point (GT)\\
\midrule
all points & 0.41 px & 0.62 px & 0.85 px\\
points seen in 3 views & 0.32 px & 0.73 px & 0.73 px\\
points seen in $\ge 9$ views & 0.55 px & 0.63 px & 1.09 px\\
\bottomrule
\end{tabular}
\end{center}
$\hat\kappa=54.6$: the tails are far heavier than Gaussian.

\section{Why not a region per detection}
For a 2D error, the region $\mathcal R=\{e\in\R^2:\ e^\top S^{-1}e\le r^2\}$ (here $e$ stands for any candidate error
vector; with $S=I$ it is the disk $e_u^2+e_v^2\le r^2$) reaches $\max_{e\in\mathcal R}e_u=r\sqrt{S_{uu}}$ along $u$ and
$r\sqrt{S_{vv}}$ along $v$ (from $\max_{e\in\mathcal R}c^\top e=r\sqrt{c^\top S c}$). It fits inside the window
$[-\epsilon,\epsilon]^2$ iff $r\le\epsilon/\sqrt{\max(S_{uu},S_{vv})}$.

The shape alone carries no probability; $\sigma$ converts it: for a Gaussian error,
$\Pr(e\notin\mathcal R)=\exp(-r^2/2\sigma^2)$. But with $M=20\,825$ detections, \eqref{eq:union} needs each region
to fail with probability $\approx\alpha/M\approx2\times10^{-6}$: that is the far tail, where the data are not Gaussian
($\hat\kappa=54.6$). A distribution-free bound does not help: Chebyshev, $\Pr(|e_u|>t)\le\sigma^2/t^2$, needs
$t\approx700\sigma$. And at equal probability the disk saves only 14\% of the square's area.
\emph{Per-detection regions therefore cannot give a useful guarantee.}

\section{Per-image ellipsoid}
For image $i$ with $n_i$ detections ($d_i=2n_i$ coordinates), define the random variable (the sum of squared errors of the image, i.e.\ the squared
Euclidean norm of its stacked error vector) and the event
\begin{equation}
  S_i=\sum_k\|e_{ik}\|^2,\qquad \mathsf E_i=\{S_i\le q_i\}.
\end{equation}
Its mean is $d_i\sigma^2$ and, if the squared coordinates are uncorrelated, its variance is $d_i\sigma^4(\kappa-1)$.
Only these two moments are needed, not the tail of a single detection.

\paragraph{Thresholds} (risk $\alpha_i$ per image, $\sum_i\alpha_i\le\alpha$):
\begin{itemize}
\item Cantelli (one-sided Chebyshev): $\Pr(S-\mu\ge t)\le s^2/(s^2+t^2)$ for mean $\mu$ and standard deviation $s$.
  Chebyshev is the two-sided version, $\Pr(|S-\mu|\ge t)\le s^2/t^2$; Cantelli is never worse for a one-sided event,
  and the two nearly coincide for small $\alpha$. Setting the bound to $\alpha_i$:
  $q_i=\hat\sigma^2\bigl(d_i+\sqrt{d_i(\hat\kappa-1)}\,\sqrt{(1-\alpha_i)/\alpha_i}\bigr)$.
\item Normal approximation (central limit theorem, $d_i$ in the hundreds):
  $q_i=\hat\sigma^2\bigl(d_i+z_{1-\alpha_i}\sqrt{d_i(\hat\kappa-1)}\bigr)$; for $\kappa=3$ this is the normal approximation of $\chi^2_{d_i}$.
\end{itemize}

\paragraph{Pose coverage.} By the proposition, $\theta^*$ is covered if \emph{some} points $X$ satisfy every
image's constraint at $\theta^*$. A witness is the least-squares point under the true cameras, so the relevant energy is
$S_i$ evaluated at that point.

\paragraph{Check on the 30-view spiral} ($\alpha=0.05$ over 30 images, $\hat\sigma=0.622$\,px and $\hat\kappa=54.6$
from leave-one-out residuals, detections flagged as mismatches removed):
\begin{center}
\begin{tabular}{lccc}
\toprule
images with $S_i\le q_i$ & normal approx. & Cantelli & Gaussian $\chi^2$\\
\midrule
pose-relevant error (least-squares point, true cameras) & 30/30 & 30/30 & 30/30\\
error to the mesh point $X^*$ & 9/30 & 30/30 & 0/30\\
\bottomrule
\end{tabular}
\end{center}
Median $S_i/(d_i\hat\sigma^2)$: 0.41 for the pose-relevant error, 2.0 for the mesh point.

\paragraph{As a convex constraint.} The depth-weighted error $v_{ik,z}\,(y_{ik}-\hat y_{ik})=A_{ik}v_{ik}$ is linear in
$(X_k,C_i)$. If $v_{ik,z}\le\bar z_{ik}$ (depth upper bound, e.g.\ from the box set), then
$\|A_{ik}v_{ik}\|/\bar z_{ik}\le\|y_{ik}-\hat y_{ik}\|$, so the truth satisfies the second-order-cone constraint
$\sum_k\|A_{ik}v_{ik}\|^2/\bar z_{ik}^{\,2}\le q_i$: an outer relaxation, no linearisation. Intersected with the box set,
the risk $\alpha$ is split between the two events.

\end{document}
